\section{The primitive-plane interface}
\label{sec:primitive-interface}

We first state the precise input inherited from TPC--70.  This section is a
dictionary, not a new estimate.

\subsection{Anchor and erasure coordinates}

Fix two distinct physical rows \(m,n\), one nonzero shift \(h_0\), a
same-residual anchor time \(j_R\), and an unequal-residual erasure time
\(j_E\).  We retain the literal hypotheses
\[
 (mn,h_0)=1,\qquad d_m\mid m,\qquad d_n\mid n,
\]
and all physical targets and residual labels used below are squarefree.
Put
\[
 H=h_0(n-m),
 \qquad
 \Delta=j_E-j_R.
\]
At the anchor, write the two targets as
\[
 A_m=mj_R+h_0=P_mR_m,
 \qquad
 A_n=nj_R+h_0=P_nR_n.
\]
At the erasure time, write
\[
 B_m=mj_E+h_0=Q_mS_m,
 \qquad
 B_n=nj_E+h_0=Q_nS_n.
\]
The four exact determinant identities are
\begin{align}
 nA_m-mA_n&=H,
 &
 nB_m-mB_n&=H,
\label{eq:two-determinants}\\
 B_m-A_m&=m\Delta,
 &
 B_n-A_n&=n\Delta,
\label{eq:vertical-differences}\\
 A_mB_n-B_mA_n&=H\Delta.
\label{eq:four-corner-determinant}
\end{align}

Suppose that the anchor residuals agree after their row factors:
\[
 d_mR_m=d_nR_n.
\]
Write
\[
 (d_m,d_n)=g,
 \qquad
 d_m=ga,\quad d_n=gb,\quad(a,b)=1.
\]
On the squarefree physical carrier, \(g,a,b\) are pairwise coprime and
squarefree.  The equality forces
\[
 R_m=bt,\qquad R_n=at
\]
for a squarefree \(t\) coprime to \(gab\), and the determinant equation
forces
\[
 t\mid |H|.
\]
For fixed \(t\), the anchor prime factors lie on
\[
 nbP_m-maP_n=H/t.
\]

At the erasure time, write
\[
 q=(S_m,S_n),
 \qquad
 S_m=qu,\quad S_n=qv,\quad(u,v)=1.
\]
Then
\[
 q\mid |H|,
 \qquad
 nuQ_m-mvQ_n=H/q.
\]
The residuals are unequal exactly when
\[
 (u,v)\ne(b,a).
\]

\subsection{Frozen sign and the aggregate kernel}

The physical target signs on the erasure edge reduce to
\begin{equation}
 \mu(d_mS_m)\mu(d_nS_n)
 =
 \mu(a)\mu(b)\mu(u)\mu(v).
\label{eq:frozen-primitive-sign}
\end{equation}
The common content \(q\) cancels, and the right side is constant along the
affine prime ray that parametrizes the admissible erasure times for fixed
\((q,u,v)\).  Thus the time variable supplies weights and masks, not a new
M\"obius oscillation.

For a fixed same-residual anchor \(\mathfrak a\), let
\[
 K_{\mathfrak a,q}(u,v)
\]
denote the complete sum along that ray, including every literal terminal
mask, endpoint, native-key constraint, prime-factor condition, normalization,
and complex amplitude, and summing only actual unequal-residual erasure
keys completing \(\mathfrak a\).  It is defined to be zero when the
corresponding ray is absent.  In particular,
\[
 K_{\mathfrak a,q}(b,a)=0.
\]
Aggregate the divisor-many contents before taking absolute values:
\begin{equation}
 K_{\mathfrak a}(u,v)
 =
 \sum_{q\mid |H|}K_{\mathfrak a,q}(u,v).
\label{eq:aggregate-q-kernel}
\end{equation}
Thus the literal aggregate is already zero at the excluded direction; set
\begin{equation}
 K^\circ_{\mathfrak a}(u,v)
 =
 K_{\mathfrak a}(u,v).
\label{eq:anchor-removed-kernel}
\end{equation}
If one instead introduces an enlarged bookkeeping kernel
\(\widehat K_{\mathfrak a}\) that includes the candidate site \((b,a)\),
the corresponding literal kernel is
\[
 K^\circ_{\mathfrak a}
 =
 \widehat K_{\mathfrak a}
 -
 \widehat K_{\mathfrak a}(b,a)e_{b,a}.
\]
That is an optional analytic convention, not an extra physical atom.
The literal signed sum attached to \(\mathfrak a\) is then
\begin{equation}
 \boxed{
 \mathfrak X_{\mathfrak a}
 =
 \mu(a)\mu(b)
 \sum_{\substack{u,v\ge1\\(u,v)=1}}
 \mu(u)\mu(v)K^\circ_{\mathfrak a}(u,v).}
\label{eq:physical-plane-coefficient}
\end{equation}
If \(\cR_{m,n}\) is the set of actual same-residual anchors for
\(m<n\), TPC--70 gives
\begin{equation}
 \operatorname{tr}(ER)
 =
 2\operatorname{Re}
 \sum_{m<n}\sum_{\mathfrak a\in\cR_{m,n}}
 \mathfrak X_{\mathfrak a}.
\label{eq:trace-by-primitive-anchors}
\end{equation}
The normalization \(W_mW_n\) is already contained in every literal
four-corner amplitude inside \(K_{\mathfrak a,q}\).

\begin{remark}[Aggregation order]
The sum over \(q\mid |H|\) in \eqref{eq:aggregate-q-kernel} precedes all
absolute values.  Bounding each \(q\)-slice separately is valid, but may pay
an unnecessary divisor factor and discard literal phase cancellation.
\end{remark}

\begin{remark}[Anchor and operator scope]
The primitive coordinates \(a,b,u,v\) depend on the row pair and the
kernel depends on the complete same-residual anchor.  A theorem for one
\(\mathfrak a\), or even cancellation after summing all anchors in
\eqref{eq:trace-by-primitive-anchors}, is not an entrywise theorem for
\(E\) and is not an operator theorem.  This distinction is retained in
\cref{sec:physical-gate}.
\end{remark}
